\documentclass[10pt]{ctexart}
\usepackage[a4paper,margin=1in]{geometry}
\usepackage{amsmath,amssymb}
\usepackage{booktabs}
\usepackage{tabularx}
\usepackage{enumitem}
\usepackage{hyperref}
\usepackage{fvextra}
\usepackage{xcolor}
\setlist[itemize]{leftmargin=1.5em}
\setlist[enumerate]{leftmargin=1.8em}
\newcommand{\code}[1]{\texttt{#1}}
\title{003：Turn-PPO 与 StepPPO 对比分析（Append-only Chat）}
\author{用户与 Codex}
\date{2026-06-08}
\begin{document}
\maketitle
\vspace{-1.2em}

\section*{记录约定}
本文件是新的 Chat 归档文件，记录本轮关于 Turn-PPO 论文、仓库实现与当前 StepPPO 的对比分析。后续如继续讨论该主题，应在本文件末尾 append，不重写已有内容。公式使用 LaTeX 形式记录。

\section*{Turn 1 | User}
\begin{Verbatim}[breaklines=true,breakanywhere=true]
有一篇Turn-PPO: Turn-Level Advantage Estimation with PPO for Improved Multi-Turn RL in Agntic LLMs，和我的想法很接近。/prj/corp/crd/morpheus/lasvegas/china-scratch/ziyuhuan/Projects/verl-agent/EXPS下下载这个论文，并查看它的git仓库。对比它和我的Step-PPO的不同之处。上边的那些分析，开一个新的/prj/corp/crd/morpheus/lasvegas/china-scratch/ziyuhuan/Projects/verl-agent/EXPS/Chat写进去
\end{Verbatim}

\section*{Turn 1 | Codex}

\subsection*{本地材料与仓库状态}
已下载和检查的材料如下：
\begin{itemize}
  \item 论文 PDF：\path{EXPS/papers/turn_ppo_2512.17008.pdf}
  \item 论文 HTML：\path{EXPS/papers/turn_ppo_2512.17008.html}
  \item 论文 HTML 文本抽取：\path{EXPS/papers/turn_ppo_2512.17008_html.txt}
  \item 原始 RAGEN 仓库浅克隆：\path{EXPS/repos/RAGEN}，remote 为 \url{https://github.com/mll-lab-nu/RAGEN.git}，当前 commit 为 \code{20daedc47558e000f7de912b060646bf2e8026bd}。
  \item 搜索到的 Turn-PPO fork：\path{EXPS/repos/RAGEN_turnPPO}，remote 为 \url{https://github.com/sjwjames/RAGEN_turnPPO.git}，当前 commit 同样为 \code{20daedc47558e000f7de912b060646bf2e8026bd}。
\end{itemize}

需要强调：论文正文没有给出一个明确的 GitHub URL，只写到训练代码基于 RAGEN 并实现了 Turn-PPO。当前搜索到的 \code{sjwjames/RAGEN\_turnPPO} 与 \code{mll-lab-nu/RAGEN} 的 main 分支在该 commit 下内容一致，\code{diff -qr} 未发现工作树差异。因此下面把“Turn-PPO 论文定义”和“公开 RAGEN 代码中可见的实现线索”分开讨论，避免把当前公开仓库误认为完整、最终的论文代码 release。

\subsection*{Turn-PPO 的核心定义}
Turn-PPO 的核心不是“每个 token 都有一个独立 MDP step”，而是把一个 agent turn 建模成 MDP 的一个 step。对第 $i$ 条 trajectory 的第 $n$ 个 turn：
\begin{itemize}
  \item state $s_{i,n}$ 是该 turn 生成前的上下文，也就是用户问题、历史交互和环境返回后的边界状态。
  \item action $a_{i,n}$ 是该 turn 的完整 LLM response，包括可能的 reasoning 和 action/tool call token：
  \[
  a_{i,n}=(a_{i,n,1},\ldots,a_{i,n,H_{i,n}}).
  \]
  \item value $V_{\phi}(s_{i,n})$ 是 turn-level state value。论文描述中，value 使用该 turn 的 query/environment-output 最后 token 对应的 critic output。
  \item advantage 在 turn 维度用 GAE 计算，而不是在 token 维度计算：
  \[
  \delta_{i,n}=r_{i,n}+\gamma V_{\phi}(s_{i,n+1})-V_{\phi}(s_{i,n}),
  \]
  \[
  A^{\mathrm{turn}}_{i,n}
  =\sum_{\ell=0}^{N_i-n}(\gamma\lambda)^\ell\delta_{i,n+\ell}.
  \]
  \item return 为
  \[
  \hat R_{i,n}=A^{\mathrm{turn}}_{i,n}+V_{\phi}(s_{i,n}).
  \]
\end{itemize}

论文的 turn-level PPO ratio 是完整 response 的 probability ratio，而不是逐 token 独立 clipping：
\[
\rho_{i,n}(\theta)
=\frac{\pi_{\theta}(a_{i,n}\mid s_{i,n})}{\pi_{\theta_{\mathrm{old}}}(a_{i,n}\mid s_{i,n})}
=\prod_{h=1}^{H_{i,n}}
\frac{\pi_{\theta}(a_{i,n,h}\mid s_{i,n},a_{i,n,<h})}
{\pi_{\theta_{\mathrm{old}}}(a_{i,n,h}\mid s_{i,n},a_{i,n,<h})}.
\]

其 actor objective 可写为：
\[
\max_{\theta}\;
\mathbb{E}\left[
\frac{1}{G}\sum_{i=1}^{G}\frac{1}{N_i}\sum_{n=1}^{N_i}
\min\left(
\rho_{i,n}(\theta)A^{\mathrm{turn}}_{i,n},
\operatorname{clip}(\rho_{i,n}(\theta),1-\epsilon,1+\epsilon)A^{\mathrm{turn}}_{i,n}
\right)
\right].
\]

critic 的 value loss 是 per-turn 的 PPO-style value objective：
\[
\min_{\phi}\;
\mathbb{E}\left[
\frac{1}{G}\sum_{i=1}^{G}\frac{1}{N_i}\sum_{n=1}^{N_i}
\frac{1}{2}\left(V_{\phi}(s_{i,n})-\hat R_{i,n}\right)^2
\right].
\]

这个定义和你最初提出的“每个 action/step 只产生一个 value，GAE/TD-error/return 在 step 维度计算”的方向非常接近。更准确地说，Turn-PPO 是 state-value 版本：value 是 $V(s_{i,n})$，不是 action-conditioned 的 $Q(s_{i,n},a_{i,n})$。因此取“生成前边界 value”比“第一个生成 token 的 value”更严格，也更符合论文定义。

\subsection*{公开 RAGEN 代码里的实现线索}
公开 RAGEN 代码中最接近 Turn-PPO 的函数是：
\begin{center}
\path{EXPS/repos/RAGEN/ragen/trainer/core_algos.py::compute_bi_level_gae_advantage_return}
\end{center}
该函数做了两层 GAE：
\begin{enumerate}
  \item high-level：在每个 turn 的 reward/EOS 位置上计算 per-turn GAE。代码注释写明 high level 是 \code{per-turn wise}。
  \item low-level：再把 high-level return 写回对应 EOS 位置，沿当前 response 的 token 位置做一遍 token-level GAE。最后对 token-level advantages 做 \code{masked\_whiten}。
\end{enumerate}

这说明公开代码仍保留了 token-level loss path 的兼容处理；它并不是一个非常干净的“只在 turn-level ratio 上 clipping”的论文公式实现。另外，\path{EXPS/repos/RAGEN/train.py} 中还有如下约束：
\[
\text{\texttt{bi\_level\_gae=True}} \Rightarrow \text{\texttt{adv\_estimator=gae}},
\]
并且当前 main 分支断言 \code{use\_turn\_scores} 在 \code{bi\_level\_gae} 下尚未正确支持。这进一步说明公开仓库不一定是论文最终 Turn-PPO 代码的完整状态。

\subsection*{当前 StepPPO-v2 的实际定义}
当前 StepPPO-v2 代码主要位于：
\begin{itemize}
  \item \path{verl/trainer/ppo/step_ppo_algos.py}
  \item \path{verl/workers/actor/step_ppo_actor.py}
  \item \path{exps/run_webshop_step_ppo_v2_4gpu_paper_align_simple.sh}
\end{itemize}

其 step-wise GAE 是：
\[
\delta^{\mathrm{step}}_{i,t}
= r^{\mathrm{step}}_{i,t}+\gamma V_{\psi}(s_{i,t+1})-V_{\psi}(s_{i,t}),
\]
\[
A^{\mathrm{step}}_{i,t}
=\delta^{\mathrm{step}}_{i,t}+\gamma\lambda A^{\mathrm{step}}_{i,t+1}.
\]
其中 trajectory 用 \code{traj\_uid} 分组，step 顺序用 \code{step\_idx} 排序。value 取 \code{pre\_action\_last\_context\_token}，也就是生成 action 前的边界 token；这点和 Turn-PPO 的 state-value 边界选择基本一致。

但 StepPPO-v2 不是纯 PPO advantage。它把 GiGPO-style episode advantage 与 step GAE advantage 混合：
\[
A^{\mathrm{final}}_{i,t}
= A^{\mathrm{epi}}_{i,t}+w\widehat A^{\mathrm{step}}_{i,t},
\quad w=0.5.
\]
当前脚本中对应配置为：
\begin{Verbatim}[breaklines=true,breakanywhere=true]
algorithm.step_ppo.step_advantage_w=0.5
algorithm.step_ppo.episode_mode=mean_norm
algorithm.step_ppo.final_advantage_mode=direct
algorithm.step_ppo.normalize_episode_advantage=False
algorithm.step_ppo.normalize_step_advantage=True
algorithm.step_ppo.normalize_final_advantage=False
actor_rollout_ref.actor.value_head.detach_value_backbone=False
actor_rollout_ref.actor.value_head.value_loss_coef=0.03
actor_rollout_ref.actor.ppo_micro_batch_size_per_gpu=8
\end{Verbatim}

这里的 critic 不是另一个单独的 model，而是 actor backbone 加一个 value MLP head。value loss 和 policy loss 在同一个 actor update path 中合并：
\[
\mathcal{L}=\mathcal{L}_{\mathrm{policy}}+c_v\mathcal{L}_{\mathrm{value}},
\quad c_v=0.03.
\]

另外，当前 StepPPO-v2 的 actor loss path 仍是 veRL 的 vanilla token-level PPO loss。即便每个 step 的 scalar advantage 被 broadcast 到该 response 的所有 token，ratio 和 clipping 仍在 token 维度做：
\[
\rho_{i,t,h}(\theta)
=\exp\left(\log\pi_\theta(a_{i,t,h})-\log\pi_{\theta_{\mathrm{old}}}(a_{i,t,h})\right).
\]
这与 Turn-PPO 论文的 response-level/turn-level ratio 不同。

\subsection*{Turn-PPO 与当前 StepPPO-v2 的主要差异}
\begin{tabularx}{\textwidth}{p{0.18\textwidth}X X}
\toprule
维度 & Turn-PPO & 当前 StepPPO-v2 \\
\midrule
MDP step & 一个 turn 是一个 MDP step，action 是完整 LLM response。 & 一个 environment/action step 是一个 step，形式上与 turn 接近。 \\
value 位置 & 生成前 state 边界；论文说取 query/environment-output 最后 token 的 critic output。 & \code{pre\_action\_last\_context\_token}，基本对齐生成前边界。 \\
advantage 来源 & 纯 critic/GAE：$A^{\mathrm{turn}}$。 & hybrid：$A^{\mathrm{epi}}+0.5\widehat A^{\mathrm{step}}$。 \\
episode term & 没有 GRPO/GiGPO episode-relative advantage 混合。 & 保留 GiGPO-style \code{mean\_norm} episode term。 \\
normalization & PPO 中通常对 GAE advantage 做 whitening；论文重点不依赖 group-relative reward normalization。 & step advantage batch-level whitening；episode advantage 用 \code{uid} group mean centering；final 不 whitening。 \\
policy ratio/clipping & response/turn-level ratio：完整 response 概率比值或其 geometric average 变体。 & 默认 \code{policy\_loss.loss\_mode=vanilla}，token-level ratio/clipping。 \\
critic 结构 & 论文说 critic 初始化自 pretrained LLM，并 attach separate value head；actor/critic 学习率分开，critic LR 约为 actor 的 5--10 倍。 & actor 自身加 value head，共享 backbone；当前无独立 critic optimizer/LR。 \\
rollout group & WebShop PPO 超参中 $G=1$，强调 PPO 更需要 prompt diversity 而非每题多 rollouts。 & 继承 WebShop StepPPO 脚本 \code{env.rollout.n=8}，仍是 grouped rollout 结构。 \\
奖励口径 & turn reward 应进入 turn-level GAE。 & step reward 已包含 invalid action penalty；但 episode term 优先读 \code{non\_tensor\_batch["episode\_rewards"]}，存在不含 invalid penalty 的风险。 \\
\bottomrule
\end{tabularx}

最关键的不同可以压缩成三点：
\begin{enumerate}
  \item Turn-PPO 是“纯 turn-level PPO/GAE”；当前 StepPPO-v2 是“GiGPO episode advantage + step GAE correction”的 hybrid。
  \item Turn-PPO 的 clipping 单位是整个 response/turn；当前 StepPPO-v2 的 clipping 单位仍是 token。
  \item Turn-PPO 使用更传统的 actor--critic 分离训练假设；当前 StepPPO-v2 使用 shared-backbone value head，把 value loss 直接加进 actor update。
\end{enumerate}

\subsection*{从 Turn-PPO 视角重新解释当前 StepPPO-v2 失败}
当前 StepPPO-v2 最近一次完整训练日志为：
\begin{center}
\path{logs/webshop_step_ppo_qwen25_15b_4gpu_paper_align_simple_20260605_031621.log}
\end{center}
该 run 完成到 step 250，没有工程崩溃，但效果明显弱于 GiGPO：
\begin{itemize}
  \item final validation：\code{val/success\_rate=0.546875}，\code{val/webshop\_task\_score=0.6619}，\code{val/text/test\_score=3.485}。
  \item best validation 大约在 step 170 附近，之后回落，说明存在 over-training 或 advantage/value signal 失真。
  \item 训练末期 \code{episode/valid\_action\_ratio} 很低，最终约 \code{0.345}，last-20 平均约 \code{0.274}；而 GiGPO 三个 seed 的 last-20 \code{valid\_action\_ratio} 约 \code{0.996--0.998}。
  \item \code{response\_length/clip\_ratio} 在 StepPPO-v2 后期明显升高，last-20 平均约 \code{0.088}；GiGPO 大约 \code{0.001} 量级。
\end{itemize}

这不是一个简单的 value loss 大小问题，而是 policy 行为已经偏离 WebShop action format：模型倾向输出长文本、中文 reasoning 或 invalid action，导致 valid action ratio 学不起来。Turn-PPO 论文本身指出，多 turn GRPO 的一个问题是 uniform trajectory advantage 不能区分不同 turn 的贡献；但 StepPPO-v2 的失败更具体：引入 step GAE 后，仍没有让 action-format 学习稳定下来。

目前最可疑的实现问题是 episode advantage 的 reward 口径。\path{verl/trainer/ppo/step_ppo_algos.py} 中 StepPPO 优先从：
\begin{center}
\code{data.non\_tensor\_batch["episode\_rewards"]}
\end{center}
读取 episode score；而 invalid action penalty 是之后写到 \code{token\_level\_scores} 和 \code{step\_rewards} 里的。也就是说：
\begin{itemize}
  \item $A^{\mathrm{step}}$ 能看到 invalid action penalty；
  \item $A^{\mathrm{epi}}$ 可能仍然使用未惩罚的 raw episode reward；
  \item v2 又采用
  \[
  A^{\mathrm{final}}=A^{\mathrm{epi}}+0.5\widehat A^{\mathrm{step}},
  \]
  因此 raw episode term 可能主导优化方向，step correction 不足以压制 invalid action。
\end{itemize}

这个问题与 Turn-PPO 的定义不一致。Turn-PPO 的 GAE reward 应该是用于训练的 turn reward；如果 invalid action penalty 是环境反馈的一部分，它必须进入 $r_{i,n}$，否则 critic 和 advantage 会学习错误的 transition value。

第二个可疑点是 final advantage scale。当前 v2 的 episode term 是 \code{mean\_norm}，其尺度可能仍接近 task reward 差值；step term 被 batch whitening 后是 unit-scale，再乘 $0.5$。如果 episode term 的数值尺度更大，实际优化会更像 GiGPO episode-only 的变体，而不是一个有效的 step-wise PPO。Turn-PPO 则不会混入一个 group-relative episode term，所以这个 scale conflict 不存在。

第三个差异是 clipping 粒度。Turn-PPO 强调完整 response 是 action，因此 PPO ratio/clipping 应该在 turn/action 维度计算。当前 StepPPO-v2 虽然 advantage 是 step scalar，但 token-level clipping 会允许同一个 action 内不同 token 受到不同 clipping 状态，严格说仍是 token-MDP loss。这可能削弱“一个 action 一个 advantage”的一致性。

第四个差异是 critic 优化。Turn-PPO 论文强调 actor/critic LR 非常敏感，并使用 actor LR $10^{-6}$、critic LR $10^{-5}$。当前 StepPPO-v2 的 value head 和 actor 共用 optimizer path；即使 \code{value\_loss\_coef=0.03}，critic 仍没有独立、更快的学习率。日志中后期 \code{actor/value\_loss} 仍经常处于较高水平，说明 value baseline 可能没有足够快地追上 step return。

\subsection*{对下一版 StepPPO 的建议}
如果目标是“更接近 Turn-PPO”，建议单独实现一个 \code{turn\_ppo} 或 \code{step\_ppo\_pure} 版本，不再继续在当前 hybrid v2 上叠补丁：
\begin{enumerate}
  \item 去掉 $A^{\mathrm{epi}}$ 混合，先做纯 step/turn GAE：
  \[
  A^{\mathrm{final}}_{i,t}=\operatorname{Whiten}(A^{\mathrm{step}}_{i,t}).
  \]
  \item 保证 $r^{\mathrm{step}}_{i,t}$ 包含 invalid action penalty、KL shaping 和所有用于 policy learning 的 reward shaping；不要让 episode term 或 return 读取未惩罚 raw reward。
  \item 将 policy loss 从 token-level clipping 切换到 response/action-level clipping。当前代码里已有 \code{compute\_policy\_loss\_gspo}，它用 response 的 geometric average ratio，可能是实现 Turn-PPO action-level clipping 的最小改动起点；但最好新建 StepPPO/TurnPPO 专用 loss，显式以每个 step response 为 action。
  \item 对 PPO WebShop 先把 group size 对齐 Turn-PPO 论文思路，尝试 \code{env.rollout.n=1}，用更高 prompt diversity 替代每题多 rollouts。GiGPO baseline 继续使用多 rollouts。
  \item critic 训练需要更强的独立性：最低限度要给 value head 单独参数组和更大学习率；更严格地则维护 separate critic model。若坚持 shared backbone，建议先试 \code{detach\_value\_backbone=True} 或 value-head-only warmup，避免 value loss 干扰 action format。
  \item 增加诊断 metrics：$A^{\mathrm{epi}}$、$A^{\mathrm{step}}$、$A^{\mathrm{final}}$ 的均值/方差/分位数，valid vs invalid action 的 advantage 均值，value explained variance，response-level ratio/clipfrac。
\end{enumerate}

如果目标是“保留 GiGPO-style hybrid StepPPO”，则至少应先修正两个问题：
\begin{enumerate}
  \item $A^{\mathrm{epi}}$ 必须从 penalized episode score 或 \code{token\_level\_rewards.sum(-1)} 计算，不能使用未含 invalid penalty 的 raw \code{episode\_rewards}。
  \item 重新标定 $A^{\mathrm{epi}}$ 与 $A^{\mathrm{step}}$ 的尺度；否则 $w=0.5$ 只是名义权重，不代表实际优化权重。
\end{enumerate}

\subsection*{结论}
Turn-PPO 与你的原始 StepPPO 想法确实很接近：二者都认为 agentic multi-turn 任务应在 action/turn/step 维度做 credit assignment，而不是把所有 token 当作等价 MDP step。当前 StepPPO-v2 已经在 value 边界和 step-level GAE 上接近 Turn-PPO，但仍有三个本质偏差：
\begin{enumerate}
  \item 它仍混入 GiGPO/GRPO episode-relative advantage；
  \item 它的 policy ratio/clipping 仍是 token-level；
  \item 它的 critic 是 shared actor value head，且没有独立 critic LR。
\end{enumerate}

因此，当前 StepPPO-v2 不能被视为 Turn-PPO 的复现。它更像一个“GiGPO episode advantage + shared-head step GAE correction”的 hybrid 方法。若要做和 Turn-PPO 更公平的比较，应该新增一个 pure Step/Turn PPO 版本，把 advantage、reward、ratio/clipping、critic optimizer 都严格对齐到 turn-level PPO 定义。

\section*{Turn 2 | User}
\begin{Verbatim}[breaklines=true,breakanywhere=true]
看一下实现的细节，value的loss计算的时候是否detach了，value head 计算的A是否detech了
\end{Verbatim}

\section*{Turn 2 | Codex}
结论：当前实现中，advantage 和 return 是 detached 的；value loss 的 target/old value 也是 detached 的；但 actor update 阶段重新计算得到的 current value prediction 默认不 detach。当前 v2 配置为 \code{detach\_value\_backbone=False}，所以 value loss 会同时反传到 value head 和 actor backbone。

\subsection*{old state values 是 detached 的}
old-log-probability 路径会同时计算 \code{old\_log\_probs} 和 \code{state\_values}。在 \path{verl/workers/actor/step_ppo_actor.py} 中，\code{compute\_log\_prob} 的 micro-batch forward 被包在：
\begin{center}
\code{with torch.no\_grad():}
\end{center}
因此，这一步生成的 \code{state\_values} 是 old value，只作为后续 GAE 的固定输入，不带梯度。

数学上，这一步得到的是：
\[
V_{\mathrm{old}}(s_t)=V_{\psi_{\mathrm{old}}}(s_t),
\]
并且在实现中等价于：
\[
\operatorname{stopgrad}\left(V_{\mathrm{old}}(s_t)\right).
\]

\subsection*{Step advantage / return 计算是 detached 的}
StepPPO 的 GAE 位于 \path{verl/trainer/ppo/step_ppo_algos.py::compute_step_gae_advantage_return}。该函数内部显式使用 \code{torch.no\_grad()}：
\begin{Verbatim}[breaklines=true,breakanywhere=true]
with torch.no_grad():
    delta = rewards[idx] + gamma * next_value * nonterminal - old_values[idx]
    lastgaelam = delta + gamma * lam * nonterminal * lastgaelam
    advantages[idx] = lastgaelam
    returns[idx] = advantages[idx] + old_values[idx]
\end{Verbatim}

所以 step-level TD-error、GAE advantage 和 return 都不会把梯度传回 value head：
\[
\delta_t
= r_t + \gamma V_{\mathrm{old}}(s_{t+1}) - V_{\mathrm{old}}(s_t),
\]
\[
A_t^{\mathrm{step}}
= \delta_t + \gamma\lambda A_{t+1}^{\mathrm{step}},
\]
\[
\hat R_t
= A_t^{\mathrm{step}} + V_{\mathrm{old}}(s_t).
\]
在实现语义上可以写成：
\[
A_t^{\mathrm{step}}
=\operatorname{stopgrad}\left(A_t^{\mathrm{step}}\right),
\quad
\hat R_t=\operatorname{stopgrad}(\hat R_t).
\]

\subsection*{Episode advantage 也是 detached 的}
GRPO-style 和 GiGPO-style episode advantage 的计算函数也都在 \code{torch.no\_grad()} 内部执行。因此：
\[
A_t^{\mathrm{epi}}
=\operatorname{stopgrad}\left(A_t^{\mathrm{epi}}\right).
\]
最终 StepPPO-v2 的 mixed advantage 为：
\[
A_t^{\mathrm{final}}
= A_t^{\mathrm{epi}} + w\widehat A_t^{\mathrm{step}},
\quad w=0.5,
\]
其中 \(A_t^{\mathrm{epi}}\) 和 \(\widehat A_t^{\mathrm{step}}\) 都是 detached training signal。代码中 \code{advantages = final\_step\_advantages.unsqueeze(-1) * response\_mask} 只是把 step scalar broadcast 到 response tokens，不会重新引入 value-head 梯度。

\subsection*{value loss 的 current value 不 detach}
actor update 阶段会重新 forward 一次 actor/value-head，得到 current values。value loss 使用：
\begin{Verbatim}[breaklines=true,breakanywhere=true]
current_values = current_values.float()
old_values = micro_data["state_values"].float()
step_returns = micro_data["step_returns"].float()

values_clipped = old_values + torch.clamp(
    current_values - old_values,
    -cliprange_value,
    cliprange_value,
)
value_loss_unclipped = (current_values - step_returns).pow(2)
value_loss_clipped = (values_clipped - step_returns).pow(2)
value_loss = 0.5 * torch.max(value_loss_unclipped, value_loss_clipped).mean()
policy_loss = policy_loss + value_loss_coef * value_loss
\end{Verbatim}

这里的梯度关系是：
\begin{itemize}
  \item \code{old\_values} 来自 old-log-prob 路径，已经 detached；
  \item \code{step\_returns} 来自 GAE 路径，已经 detached；
  \item \code{current\_values} 来自当前 actor update forward，不 detach；
  \item 因此 \code{value\_loss} 会对 \code{current\_values} 反传。
\end{itemize}

对应公式为：
\[
\mathcal{L}_{V}
=\frac{1}{2}\max\left(
\left(V_{\psi}(s_t)-\operatorname{stopgrad}(\hat R_t)\right)^2,
\left(\operatorname{clip}\left(V_{\psi}(s_t),V_{\mathrm{old}}(s_t)-\epsilon_v,V_{\mathrm{old}}(s_t)+\epsilon_v\right)-\operatorname{stopgrad}(\hat R_t)\right)^2
\right).
\]
因此：
\[
\nabla_{\psi}\mathcal{L}_{V}\neq 0.
\]

\subsection*{是否反传到 backbone 由 detach\_value\_backbone 控制}
value head 的 forward 在 \path{verl/workers/actor/value_head.py} 中实现：
\begin{Verbatim}[breaklines=true,breakanywhere=true]
hidden_states = output.hidden_states[-1]
if bool(getattr(self, "_step_value_head_detach_backbone", False)):
    hidden_states = hidden_states.detach()
step_values = self.value_head(hidden_states)
\end{Verbatim}

因此：
\begin{itemize}
  \item 若 \code{detach\_value\_backbone=True}，value loss 只更新 value head，不通过 value path 更新 actor backbone；
  \item 若 \code{detach\_value\_backbone=False}，value loss 同时更新 value head 和 actor backbone。
\end{itemize}

当前 v2 脚本 \path{exps/run_webshop_step_ppo_v2_4gpu_paper_align_simple.sh} 中设置：
\begin{center}
\code{actor\_rollout\_ref.actor.value\_head.detach\_value\_backbone=False}
\end{center}
所以当前 v2 的实际梯度为：
\[
\nabla_{\theta_{\mathrm{backbone}}}\mathcal{L}_{V}\neq 0,
\quad
\nabla_{\psi_{\mathrm{value\_head}}}\mathcal{L}_{V}\neq 0.
\]
如果改为 \code{detach\_value\_backbone=True}，则变成：
\[
\nabla_{\theta_{\mathrm{backbone}}}\mathcal{L}_{V}=0,
\quad
\nabla_{\psi_{\mathrm{value\_head}}}\mathcal{L}_{V}\neq 0.
\]
policy loss 仍然会正常更新 actor backbone。

\subsection*{一句话总结}
当前 StepPPO 实现是 PPO-style 的 stop-gradient 结构：\(A\) 和 \(\hat R\) 都是 old policy/value 下的 detached 训练信号；value loss 的 \(V_{\psi}(s_t)\) 不 detach。由于当前 v2 使用 \code{detach\_value\_backbone=False}，value loss 会反传到 shared actor backbone，这可能干扰 policy/action-format 学习，也是后续 ablation 中应重点检查的因素。

\section*{Turn 3 | User}
\begin{Verbatim}[breaklines=true,breakanywhere=true]
2. 保证 r
step
i,t
包含 invalid action penalty、KL shaping 和所有用于 policy learning 的 reward shaping；
不要让 episode term 或 return 读取未惩罚 raw reward。这里的reward shaping具体是？写到003中
\end{Verbatim}

\section*{Turn 3 | Codex}
这里的 reward shaping 指的是：在环境原始奖励 \(r^{\mathrm{env}}_{i,t}\) 之外，为了让 policy learning 感知格式、约束、偏好或稳定性而加入到 reward 空间里的所有修正项。它们不是独立的 logging metric，而是会进入 advantage/return 计算的训练奖励。对 StepPPO 来说，严格要求是：critic/GAE 使用的 \(r^{\mathrm{step}}_{i,t}\)、episode term 使用的 score，以及 policy loss 间接受到的 return/advantage，必须基于同一套 shaped reward 口径。

\subsection*{一般形式}
可以把 step reward 写成：
\[
r^{\mathrm{step}}_{i,t}
= r^{\mathrm{env}}_{i,t}
+ r^{\mathrm{task\_shape}}_{i,t}
- \alpha\mathbf{1}\{a_{i,t}\;\text{invalid}\}
- \beta\sum_{h\in a_{i,t}}D_{\mathrm{KL},i,t,h}
+ r^{\mathrm{other\_shape}}_{i,t}.
\]
其中：
\begin{itemize}
  \item \(r^{\mathrm{env}}_{i,t}\)：环境返回的原始 step reward，例如 WebShop 环境在某一步返回的分数增量或最终得分信号。
  \item \(r^{\mathrm{task\_shape}}_{i,t}\)：任务侧 dense shaping，例如 step bonus、process reward、progress reward、长度归一化后的 reward 等。如果当前任务没有这些项，则为 0。
  \item \(-\alpha\mathbf{1}\{a_{i,t}\;\text{invalid}\}\)：invalid action penalty，用于惩罚格式错误、非法 tool call 或环境无法执行的 action。
  \item \(-\beta\sum_h D_{\mathrm{KL},i,t,h}\)：reward-space KL penalty。只有当 KL 被配置为 reward shaping，即 \code{algorithm.use\_kl\_in\_reward=True} 时，它才属于 \(r^{\mathrm{step}}\)。
  \item \(r^{\mathrm{other\_shape}}_{i,t}\)：其他真正进入 reward 的修正项，例如 safety penalty、重复动作 penalty、action length penalty、tool-call cost 等。
\end{itemize}

这里的原则是：只要某个项会改变 advantage/return，就必须被视为 reward shaping；只用于日志或只作为额外 loss 的项，不属于 \(r^{\mathrm{step}}\)。

\subsection*{当前代码里的实际 reward 口径}
当前 rollout 中，环境每一步返回 \code{rewards}，并写入：
\begin{center}
\code{batch.non\_tensor\_batch["rewards"]}
\end{center}
同时累计得到 raw episode reward：
\begin{center}
\code{episode\_rewards += rewards}
\end{center}
reward manager 后续会把 \code{episode\_rewards} 写到 \code{token\_level\_scores} 的最后一个 response token 上。也就是说，raw episode score 的来源是环境累计分数。

随后训练 loop 会先应用 invalid action penalty。当前实现会同时修改 token-level score 和 step reward：
\begin{Verbatim}[breaklines=true,breakanywhere=true]
reward_tensor[i, valid_response_length - 1] -= invalid_action_penalty_coef * action_invalids
step_rewards[i] -= invalid_action_penalty_coef * action_invalids
\end{Verbatim}
因此，在当前配置下，invalid action penalty 是真正的 reward shaping。它应该进入：
\[
r^{\mathrm{step}}_{i,t}
= r^{\mathrm{env}}_{i,t}-0.1\cdot\mathbf{1}\{a_{i,t}\;\text{invalid}\},
\]
因为当前脚本里 \code{invalid\_action\_penalty\_coef=0.1}。

\subsection*{KL shaping 和 actor KL loss 的区别}
当前代码有两种 KL 使用方式，必须区分：
\begin{enumerate}
  \item reward-space KL penalty：当 \code{algorithm.use\_kl\_in\_reward=True} 时，代码执行
  \begin{center}
  \code{token\_level\_rewards = token\_level\_scores - beta * kld}
  \end{center}
  这属于 reward shaping，因为它改变了用于 advantage 的 reward。
  \item actor KL loss：当 \code{actor\_rollout\_ref.actor.use\_kl\_loss=True} 时，KL 作为 policy loss 的额外正则项加入 actor update。这不是 reward shaping，因为它不改变 \(r^{\mathrm{step}}\)、return 或 advantage。
\end{enumerate}

当前 StepPPO-v2 的 WebShop 脚本是：
\begin{center}
\code{algorithm.use\_kl\_in\_reward=False}
\end{center}
并继承：
\begin{center}
\code{actor\_rollout\_ref.actor.use\_kl\_loss=True}
\end{center}
所以当前 v2 中 KL 是 actor loss 正则，不是 \(r^{\mathrm{step}}\) 的 reward shaping。当前 v2 的 \(r^{\mathrm{step}}\) 至少应该包含环境 step reward 和 invalid action penalty；不需要包含 actor KL loss。

如果后续打开 \code{algorithm.use\_kl\_in\_reward=True}，则 StepPPO 需要额外小心：现有 \code{apply\_kl\_penalty} 主要更新 \code{token\_level\_rewards}，但 StepPPO 的 GAE 使用的是 \code{step\_rewards}。因此，若 KL penalty 被放入 reward 空间，就必须把每个 step response token 上的 KL penalty 聚合回对应 \code{step\_rewards}：
\[
r^{\mathrm{step}}_{i,t}
= r^{\mathrm{env}}_{i,t}
-\alpha\mathbf{1}\{a_{i,t}\;\text{invalid}\}
-\beta\sum_{h\in a_{i,t}}D_{\mathrm{KL},i,t,h}.
\]
否则会出现 token-level episode advantage 看到了 KL penalty，而 step GAE return 没看到 KL penalty 的不一致。

\subsection*{为什么不能读取未惩罚 raw episode reward}
当前最需要避免的是：\(A^{\mathrm{step}}\) 使用 shaped step reward，但 \(A^{\mathrm{epi}}\) 仍从未惩罚的 raw \code{episode\_rewards} 读取。这样会导致：
\[
A^{\mathrm{final}}_{i,t}
= A^{\mathrm{epi,raw}}_{i,t}+wA^{\mathrm{step,shaped}}_{i,t},
\]
其中 episode term 可能奖励一个最终得分不错但中间 invalid action 很多的 trajectory，而 step term 试图惩罚 invalid action。两个信号方向冲突时，final advantage 会变得不可靠。

更一致的做法是定义 shaped episode score：
\[
R^{\mathrm{epi,shaped}}_i
=\sum_{t=1}^{T_i} r^{\mathrm{step,shaped}}_{i,t},
\]
然后 episode advantage 也从 \(R^{\mathrm{epi,shaped}}_i\) 计算，而不是从 raw \code{episode\_rewards} 计算。

\subsection*{当前 StepPPO 的具体含义}
对当前 v2 配置，reward shaping 具体指：
\begin{itemize}
  \item 必须包含：环境 step reward \(r^{\mathrm{env}}_{i,t}\)。
  \item 必须包含：invalid action penalty，即 \(-0.1\cdot\mathbf{1}\{a_{i,t}\;\text{invalid}\}\)。
  \item 当前不包含：reward-space KL penalty，因为 \code{algorithm.use\_kl\_in\_reward=False}。
  \item 当前不包含：actor KL loss，因为它是额外 objective，不进入 reward/return。
  \item 如果未来加入 process reward、step bonus、tool cost、length penalty、safety penalty 等，只要它们用于改变 advantage/return，就也必须进入 \(r^{\mathrm{step}}\) 和 shaped episode score。
\end{itemize}

一句话总结：这里的 reward shaping 不是泛泛地“奖励工程”，而是指所有进入 \(r^{\mathrm{step}}\)、return、advantage 的 reward-space 修正。当前 v2 最重要的是保证 invalid action penalty 同时进入 step GAE 和 episode advantage；KL 只有在 \code{use\_kl\_in\_reward=True} 时才需要进入 \(r^{\mathrm{step}}\)，否则它只是 actor KL loss。

\section*{Turn 4 | User}
\begin{Verbatim}[breaklines=true,breakanywhere=true]
和GiGPO对齐吧，A_epi的做法。然后A_step和GiGPO的step adv同样做法，这样如何
\end{Verbatim}

\section*{Turn 4 | Codex}
可以，但需要先明确这个实验的算法性质：如果 \(A^{\mathrm{epi}}\) 和 \(A^{\mathrm{step}}\) 都完全按 GiGPO 做，那么核心 advantage 已经是 GiGPO-style relative advantage；剩下的差异只来自是否额外保留 value head/value loss，以及 actor loss path 是否仍与 GiGPO 完全一致。此时它不再是严格意义上的 value-based Step-PPO/Turn-PPO。

对齐后的 advantage 形式是：
\[
A_{i,t}^{\mathrm{final}}
=
A_{i,t}^{\mathrm{epi,GiGPO}}
+wA_{i,t}^{\mathrm{step,GiGPO}}.
\]
其中 episode term 为：
\[
A_{i,t}^{\mathrm{epi,GiGPO}}
=
\operatorname{Norm}_{\mathrm{uid}}\left(R_i^{\mathrm{epi,shaped}}\right),
\]
step term 为：
\[
A_{i,t}^{\mathrm{step,GiGPO}}
=
\operatorname{Norm}_{\mathrm{anchor}(s_{i,t}),\mathrm{uid}}
\left(G_{i,t}^{\mathrm{step,shaped}}\right).
\]
这里的 \(G_{i,t}^{\mathrm{step,shaped}}\) 是 GiGPO 的 step discounted return，而不是 critic GAE：
\[
G_{i,t}^{\mathrm{step,shaped}}
=
\sum_{k=t}^{T_i}\gamma^{k-t}r_{i,k}^{\mathrm{step,shaped}}.
\]
这意味着不再使用 critic TD-error：
\[
\delta_{i,t}=r_{i,t}+\gamma V(s_{i,t+1})-V(s_{i,t}),
\]
也不再使用 value-based GAE：
\[
A_{i,t}^{\mathrm{GAE}}
=
\sum_{l\ge 0}(\gamma\lambda)^l\delta_{i,t+l}.
\]
因此，这个版本更准确地应称为 GiGPO-aligned advantage 诊断实验，而不是 Step-PPO 主算法。

\subsection*{这样做的价值}
这个实验很有诊断价值，因为它可以排除当前 StepPPO-v2 中两个主要不一致点：
\begin{enumerate}
  \item \(A^{\mathrm{epi}}\) 的 reward 口径不一致。当前 StepPPO-v2 可能从 raw \code{episode\_rewards} 读取 episode score，而 GiGPO 从 shaped \code{token\_level\_rewards.sum(-1)} 计算 episode advantage。
  \item \(A^{\mathrm{step}}\) 的定义和归一化不一致。当前 StepPPO-v2 用 value-based GAE 后做 batch-level whitening；GiGPO 用 step discounted return，并在 \code{anchor\_obs + uid} 定义的 step group 内做 relative normalization。
\end{enumerate}

如果这个 GiGPO-aligned 版本恢复到 GiGPO 的表现，说明当前 v2 的失败主要来自 value-based GAE、value loss 或 reward 口径不一致。如果它仍然不稳定，则需要继续检查 actor update path、token-level clipping、value auxiliary loss、micro-batch 或其它工程差异。

\subsection*{是否保留 value head}
需要拆开两个版本，否则实验含义会混在一起。

第一个版本是 \textbf{GiGPO-aligned-no-value}：
\begin{itemize}
  \item \(A^{\mathrm{epi}}\) 完全使用 GiGPO 做法；
  \item \(A^{\mathrm{step}}\) 完全使用 GiGPO 做法；
  \item 关闭 value head 和 value loss；
  \item 这个版本应当接近 GiGPO baseline，用于验证对齐是否正确。
\end{itemize}

第二个版本是 \textbf{GiGPO-aligned-value-aux}：
\begin{itemize}
  \item advantage 仍完全使用 GiGPO；
  \item 保留 value head，但 value head 只作为 auxiliary value regression；
  \item critic 不参与 advantage；
  \item 建议先设置 \code{detach\_value\_backbone=True}，避免 auxiliary value loss 干扰 actor backbone；
  \item \code{value\_loss\_coef} 保持较小，例如 \code{0.01} 或 \code{0.03}。
\end{itemize}

如果保留 value head 且 \code{detach\_value\_backbone=False}，则算法变成：
\[
\mathcal{L}
=
\mathcal{L}_{\mathrm{GiGPO\;policy}}
+c_v\mathcal{L}_{V},
\]
且 \(\mathcal{L}_{V}\) 会反传到 shared actor backbone。这样即使 advantage 已经对齐 GiGPO，value loss 仍可能改变 actor 表征并影响 action format 学习。因此不建议把这个混合版本作为第一个诊断实验。

\subsection*{与真正 Step-PPO/Turn-PPO 的关系}
如果目标是研究 Step-PPO/Turn-PPO，那么最终主线仍应使用 value-based GAE：
\[
A_{i,t}^{\mathrm{step}}
=
\sum_{l\ge 0}(\gamma\lambda)^l
\left(
 r_{i,t+l}
 +\gamma V(s_{i,t+l+1})
 -V(s_{i,t+l})
\right).
\]
GiGPO-aligned 版本不能证明 critic-based Step-PPO 有效；它只能作为 diagnostic control，用来定位当前失败到底来自 reward/normalization 不一致，还是来自 value head/value loss 本身。

\subsection*{建议命名}
建议新增两个诊断脚本：
\begin{itemize}
  \item \code{v3\_gigpo\_aligned\_no\_value}：完全 GiGPO advantage，无 value head。
  \item \code{v3\_gigpo\_aligned\_value\_aux}：完全 GiGPO advantage，加 auxiliary value loss。
\end{itemize}

其中 \code{v3\_gigpo\_aligned\_no\_value} 应作为第一优先级。如果它不能复现 GiGPO 附近的结果，说明实现还没有真正对齐；如果它能复现，再加入 \code{value\_aux} 检查 value head 是否会破坏稳定性。

\section*{Turn 5 | User}
\begin{Verbatim}[breaklines=true,breakanywhere=true]
target应该是StepPPO的那个GAE return。为什么不这么做

两套方案都先说一下怎么做吧，先记录到003中
\end{Verbatim}

\section*{Turn 5 | Codex}
这里需要区分两套 value-aux target 方案。二者都保持 GiGPO policy advantage 不变：
\[
A_{i,t}^{\mathrm{policy}}
=
A_{i,t}^{\mathrm{epi,GiGPO}}
+wA_{i,t}^{\mathrm{step,GiGPO}}.
\]
区别只在 value head 回归什么 target。

\subsection*{方案 A：回归 GiGPO step return}
这是当前 \code{gigpo\_v1\_value\_aux} 里最小改动的实现方式。GiGPO 原本会先计算 step discounted return：
\[
G_{i,t}^{\mathrm{GiGPO}}
=
\sum_{k=t}^{T_i}\gamma^{k-t}r_{i,k}.
\]
然后 GiGPO 用它做 step relative advantage：
\[
A_{i,t}^{\mathrm{step,GiGPO}}
=
\operatorname{Norm}_{\mathrm{anchor}(s_{i,t}),\mathrm{uid}}
\left(G_{i,t}^{\mathrm{GiGPO}}\right).
\]
方案 A 直接把这个 GiGPO step return 作为 value target：
\[
Y_{i,t}^{A}=G_{i,t}^{\mathrm{GiGPO}}.
\]
实现上，在 GiGPO advantage 计算完成后补：
\begin{center}
\code{data.batch["step\_returns"] = data.batch["step\_rewards"].detach().clone()}
\end{center}
然后复用 shared value-head actor path 里的 clipped value loss：
\[
\mathcal{L}_{V}^{A}
=\frac{1}{2}\max\left(
\left(V_{\psi}(s_{i,t})-\operatorname{stopgrad}(Y_{i,t}^{A})\right)^2,
\left(\operatorname{clip}\left(V_{\psi}(s_{i,t}),V_{\mathrm{old}}(s_{i,t})-\epsilon_v,V_{\mathrm{old}}(s_{i,t})+\epsilon_v\right)-\operatorname{stopgrad}(Y_{i,t}^{A})\right)^2
\right).
\]

这个方案的优点是实现非常小，不改变 GiGPO advantage，也不需要重建 step order 或 GAE。它测试的是：value head 是否能作为 auxiliary task 回归 GiGPO 已经使用的 step signal，以及这个 auxiliary loss 是否会破坏 GiGPO 稳定性。

但它的限制也很明确：
\begin{itemize}
  \item 它不是 PPO/StepPPO critic target，因为没有使用 bootstrap value，也没有使用 TD-error/GAE。
  \item 它更像 Monte-Carlo step return regression，critic 不参与 advantage。
  \item 当前 GiGPO pipeline 是先用 raw immediate reward 计算 discounted return，再应用 invalid action penalty 到 \code{step\_rewards} 当前 row；严格来说，未来 step 的 invalid penalty 没有折扣传播回更早的 return。
  \item 如果目标是判断 value head 能否成为真正 critic，方案 A 的证据较弱。
\end{itemize}

因此，方案 A 是一个低风险诊断：\textbf{GiGPO policy + auxiliary regression to GiGPO step signal}。它可以回答“加 value head/value loss 是否会破坏 GiGPO”，但不能充分回答“StepPPO critic 是否可行”。

\subsection*{方案 B：回归 StepPPO GAE return}
这是你指出的更合理的 critic-oriented 方案。policy advantage 仍然保持 GiGPO 不变，但 value head 的 target 改为 StepPPO 那套 GAE return：
\[
\delta_{i,t}
=
 r_{i,t}^{\mathrm{step,shaped}}
 +\gamma V_{\mathrm{old}}(s_{i,t+1})
 -V_{\mathrm{old}}(s_{i,t}),
\]
\[
A_{i,t}^{\mathrm{GAE}}
=
\sum_{l\ge 0}(\gamma\lambda)^l\delta_{i,t+l},
\]
\[
Y_{i,t}^{B}
=
\hat R_{i,t}^{\mathrm{GAE}}
=
A_{i,t}^{\mathrm{GAE}}+V_{\mathrm{old}}(s_{i,t}).
\]
value loss 变成：
\[
\mathcal{L}_{V}^{B}
=\frac{1}{2}\max\left(
\left(V_{\psi}(s_{i,t})-\operatorname{stopgrad}(\hat R_{i,t}^{\mathrm{GAE}})\right)^2,
\left(\operatorname{clip}\left(V_{\psi}(s_{i,t}),V_{\mathrm{old}}(s_{i,t})-\epsilon_v,V_{\mathrm{old}}(s_{i,t})+\epsilon_v\right)-\operatorname{stopgrad}(\hat R_{i,t}^{\mathrm{GAE}})\right)^2
\right).
\]

实现上需要做几件事：
\begin{enumerate}
  \item old-log-probability 阶段仍然让 value head 输出 \code{state\_values}，这个值在 \code{torch.no\_grad()} 下计算，是 \(V_{\mathrm{old}}(s_{i,t})\)。
  \item 在 batch 被 \code{adjust\_batch} 复制前，为每个真实 step 记录 \code{step\_idx} 和 duplicate-safe 的 \code{step\_sample\_uid}，避免 padding/复制样本被 GAE 当成额外 environment transition。
  \item 为 value target 单独构造 shaped immediate reward \(r_{i,t}^{\mathrm{step,shaped}}\)。不能直接用 GiGPO 的 \code{step\_rewards}，因为那已经是 discounted return，而不是 immediate reward。
  \item 当前配置下 shaped immediate reward 至少应为
  \[
  r_{i,t}^{\mathrm{step,shaped}}
  =r_{i,t}^{\mathrm{env}}-0.1\cdot\mathbf{1}\{a_{i,t}\;\text{invalid}\}.
  \]
  如果未来打开 \code{algorithm.use\_kl\_in\_reward=True}，还要把 step response token 上的 KL penalty 聚合进 \(r_{i,t}^{\mathrm{step,shaped}}\)。
  \item 用 StepPPO 的 GAE helper 计算 \(A^{\mathrm{GAE}}\) 和 \(\hat R^{\mathrm{GAE}}\)，但只把 \(\hat R^{\mathrm{GAE}}\) 写入 \code{step\_returns} 作为 value loss target；policy loss 的 \code{advantages} 仍然保留 GiGPO 的 \(A^{\mathrm{policy}}\)。
\end{enumerate}

这个方案的优点是：
\begin{itemize}
  \item value head 的训练目标和 StepPPO/Turn-PPO critic 语义一致；
  \item 能直接测试 value head 是否能学习 bootstrapped step return；
  \item 如果这个 auxiliary target 稳定，而 StepPPO-v2 不稳定，就说明问题更可能出在“用 GAE advantage 更新 policy”而不只是 value head 本身。
\end{itemize}

它的限制是：
\begin{itemize}
  \item 实现比方案 A 更侵入，需要额外维护 immediate shaped reward、step order 和 duplicate-safe id；
  \item policy advantage 与 value target 不再来自同一个 estimator：policy 用 GiGPO relative advantage，value 用 GAE return，所以它仍然不是标准 PPO，只是更接近 critic feasibility test；
  \item 如果 \code{detach\_value\_backbone=False}，value loss 仍可能通过 shared backbone 影响 GiGPO policy 表征。
\end{itemize}

\subsection*{为什么一开始没有直接做方案 B}
方案 A 是最小诊断补丁：复用 GiGPO 已存在的 \code{step\_rewards}，只验证“加 value head/value loss 会不会破坏 GiGPO”。但你指出的是对的：如果我们真正想测试 StepPPO critic 是否可行，target 应该是方案 B 的 \(\hat R^{\mathrm{GAE}}\)，而不是 GiGPO step return。

因此后续更合理的实验顺序是：
\begin{enumerate}
  \item 先跑方案 A，回答最小问题：auxiliary value loss 是否破坏 GiGPO 稳定性。
  \item 再跑方案 B，回答更关键的问题：shared value head 能否学习 StepPPO/Turn-PPO 风格的 GAE return。
  \item 若方案 B 稳定，再考虑让 \(A^{\mathrm{GAE}}\) 进入 policy advantage；否则不要直接把 GAE advantage 加进 policy。
\end{enumerate}

命名上建议：
\begin{itemize}
  \item \code{gigpo\_v1\_value\_aux\_return\_gigpo}：方案 A。
  \item \code{gigpo\_v1\_value\_aux\_return\_gae}：方案 B。
\end{itemize}

一句话结论：方案 A 适合测试 value loss 是否破坏 GiGPO；方案 B 才适合测试 StepPPO critic target 是否可行。若只能选一个与 StepPPO 目标更一致的版本，应优先实现方案 B。

\section*{Turn 6 | User}
\begin{Verbatim}[breaklines=true,breakanywhere=true]
你应该使用的是方案 B这种实现方式
\end{Verbatim}

\section*{Turn 6 | Codex}
已将 \code{gigpo\_v1\_value\_aux} 的实现口径修正为方案 B：policy advantage 仍然完全使用 GiGPO，但 auxiliary value loss 的 target 改为 StepPPO-style GAE return，而不是 GiGPO step discounted return。

修正后的实现语义为：
\[
A_{i,t}^{\mathrm{policy}}
=
A_{i,t}^{\mathrm{epi,GiGPO}}
+wA_{i,t}^{\mathrm{step,GiGPO}},
\]
但 value head target 为：
\[
\delta_{i,t}
=
 r_{i,t}^{\mathrm{step,shaped}}
 +\gamma V_{\mathrm{old}}(s_{i,t+1})
 -V_{\mathrm{old}}(s_{i,t}),
\]
\[
A_{i,t}^{\mathrm{GAE}}
=
\sum_{l\ge 0}(\gamma\lambda)^l\delta_{i,t+l},
\]
\[
\hat R_{i,t}^{\mathrm{GAE}}
=A_{i,t}^{\mathrm{GAE}}+V_{\mathrm{old}}(s_{i,t}).
\]
然后写入：
\begin{center}
\code{data.batch["step\_returns"] = stopgrad(}\(\hat R_{i,t}^{\mathrm{GAE}}\)\code{)}
\end{center}
触发 shared value-head path 中的 clipped value loss。

具体实现点如下：
\begin{itemize}
  \item 新增 trainer 仍为 \path{verl/trainer/ppo/gigpo_v1_value_aux_trainer.py}。
  \item 原始 GiGPO advantage 不变，仍由原始 \code{compute\_gigpo\_outcome\_advantage} 生成。
  \item 在 \code{adjust\_batch} 前补 \code{step\_idx} 和 \code{step\_sample\_uid}，避免 batch padding/copy 产生伪 transition。
  \item value target 单独从 raw per-step \code{non\_tensor\_batch["rewards"]} 构造 shaped immediate reward，而不是复用 GiGPO 的 discounted \code{step\_rewards}。
  \item 当前 shaped immediate reward 至少包含 invalid action penalty：
  \[
  r_{i,t}^{\mathrm{step,shaped}}
  =r_{i,t}^{\mathrm{env}}-\alpha\mathbf{1}\{a_{i,t}\;\mathrm{invalid}\}.
  \]
  \item 如果未来打开 \code{algorithm.use\_kl\_in\_reward=True}，实现会把 \code{token\_level\_rewards - token\_level\_scores} 的 token sum 作为 reward-space KL delta 加回当前 step immediate reward。
  \item 配置已改为：
  \begin{center}
  \code{algorithm.gigpo\_v1\_value\_aux.value\_target = "gae\_returns"}
  \end{center}
\end{itemize}

因此，当前 \code{gigpo\_v1\_value\_aux} 不再是方案 A，而是：
\[
\mathcal{L}
=
\mathcal{L}_{\mathrm{GiGPO\;policy}}
+c_v\mathcal{L}_{V}^{\mathrm{GAE-return}}.
\]
它仍不是标准 PPO，因为 policy advantage 仍为 GiGPO；但它更适合测试 shared value head 是否能学习 StepPPO/Turn-PPO 风格的 critic target。


\section*{Turn 7 | User}
\begin{Verbatim}[breaklines=true,breakanywhere=true]
接下来我有个想法，GiGPO是使用硬文本匹配来计算step的advatage的。有没有一种其他方法？比如模型的既视感之类的，来对各个step之间进行软匹配，从而实现估计。这个的好处是不需要硬性匹配了。

请你草拟这个想法为一个文档
\end{Verbatim}

\section*{Turn 7 | Codex}
Codex 将这个想法整理为一个新的 GiGPO 扩展方向，暂命名为 \textbf{Soft-GiGPO / Semantic Step Matching}。核心观点是：当前 GiGPO 的 step group 默认依赖 exact anchor observation match，可选 \code{SequenceMatcher} 也仍然是阈值式 hard clustering；新的设想把“相同 state”的硬文本定义放宽成“相似 state”的连续语义定义。

文档中定义了 soft matching 的基本形式。给定 step 表示 $z_i=f_{\mathrm{match}}(s_i)$ 和相似度 $\mathrm{sim}_{ij}=\cos(z_i,z_j)$，在同一 task \code{uid} 内构造 leave-one-out softmax 权重：
\[
w_{ij}=\frac{\mathbf{1}[u_i=u_j]\mathbf{1}[i\ne j]\exp(\mathrm{sim}_{ij}/\tau)}
{\sum_{k:u_k=u_i,k\ne i}\exp(\mathrm{sim}_{ik}/\tau)+\epsilon}.
\]
再用加权 baseline
\[
b_i^{\mathrm{soft}}=\sum_j w_{ij}R_j
\]
得到 soft step advantage：
\[
A_i^{\mathrm{soft}}=R_i-b_i^{\mathrm{soft}}.
\]

文档建议第一版只做 \textbf{Soft-GiGPO-v0}：不改变 policy loss、不接入 StepPPO-v3 value head、不改变 rollout recipe，只替换或增强 GiGPO 的 step relative advantage estimator。候选匹配表示包括 text embedding、actor selected hidden state、frozen actor snapshot，以及后续更高风险的 value-aware similarity。

新增文档位置：
\begin{itemize}
  \item \path{EXPS/019_soft_gigpo_semantic_step_matching.md}
  \item \path{EXPS/tex/019_soft_gigpo_semantic_step_matching.tex}
  \item \path{EXPS/019_soft_gigpo_semantic_step_matching.pdf}
\end{itemize}

\end{document}
